\scanpage{549}

and we assume that $\mathcal{G}^+\ne\mathcal{G}$, hence
\[
\tag{3}
\begin{gathered}
\mathcal{G}\simeq\{1,\tau'\}\mathbin{.}\mathcal{G}^+\\[-2pt]
(\tfrac12\ \text{direct})
\end{gathered}
\]
and we have a canonical exact sequence
\[
\tag{4}
\diagram{\begin{tikzcd}[column sep=large]
1\arrow[r]&\mathcal{G}^+\arrow[r]&\mathcal{G}\arrow[r,"\varepsilon_{\mathcal{G}}"]&\{\pm1\}\arrow[r]&1
\end{tikzcd}}
\]
We also denote by $\nu_{\mathcal{G}}$ the composite
\[
\tag{5}
\diagram{\begin{tikzcd}[column sep=large,row sep=small]
\mathcal{G}\arrow[r,"\varepsilon_{\mathcal{G}}"]\arrow[rr,bend left=25,"\nu_{\mathcal{G}}"]&\{\pm1\}&\mathbb{Z}/2\mathbb{Z}\arrow[l,"\sim"']\\
&(-1)^\nu&\nu\arrow[l,mapsto]
\end{tikzcd}}\ .
\]

In general, we \unsure{will denote} (by abuse of language) by the same letter elements of $\mathrm{Gl}(2,\mathbb{Z})$ --- such as $\varepsilon_0,\varepsilon_1,\tau=\tau_\infty,\tau'=\tau'_\infty$, $\sigma=\sigma_\infty,\rho=\rho_\infty$ etc --- and their images in $\mathcal{G}$.

\margin{The groups $Z_{\rho,\sigma}$, $\gamma_{\rho,\sigma}$}

$2^\circ)$ Let
\[
\tag{6}
\begin{gathered}
\deleted{B}Z_{\rho,\sigma}\ \subset\ \operatorname{Aut}(\mathcal{G}^+)\times\widehat{\mathbb{Z}}^*\times\{\pm1\}\times\{\pm1\}\\[-2pt]
\Vert\ \mathrm{dfn}\\[-2pt]
\left\{(u,\mu,\varepsilon,\varepsilon')\ \middle|\
\begin{aligned}
u(\varepsilon_0)&=\varepsilon_0^\mu\\
u(\rho)&\ \text{conjugate in }\mathcal{G}^+\text{ to }\rho^{\varepsilon'}\\
u(\sigma)&\ \rule{4.1cm}{0.3pt}\ \text{to }\sigma^\varepsilon
\end{aligned}\right\}
\end{gathered}
\]
This is a subgroup of the product group.

Let us put
\[
\tag{7}
\left\{\begin{aligned}
L_\rho&=\text{subgroup of }\mathcal{G}\text{ generated by }\rho\\
L_\sigma&=\rule{7cm}{0.3pt}\ \sigma
\end{aligned}\right.
\]
so that we have \unsure{epimorphisms}
\[
\tag{8}
\left\{\diagram{\begin{tikzcd}[column sep=large,row sep=small]
\mathbb{Z}/6\mathbb{Z}\arrow[r,"\text{epimorphism}"]&L_\rho&n\arrow[r,mapsto]&\rho^n\\
\mathbb{Z}/4\mathbb{Z}\arrow[r]&L_\sigma&n\arrow[r,mapsto]&\unsure{\sigma^n}
\end{tikzcd}}\right.\ ,
\]
\mspage{548}

{\small
\begin{center}
548\\
\texttt{https://grothendieck.umontpellier.fr}\\
Université de Montpellier, UMR CNRS 5149, Case courrier 051, Place Eugène Bataillon\\
34095 Montpellier cedex 5 - France
\end{center}
}

\scanpage{550}
\mspage{578}

and \deleted{some} \inserted{a commutative\ diagram of }homomorphisms
\[
\tag{9}
\diagram{\begin{tikzcd}[column sep=large,row sep=small]
&L_\rho\arrow[dr,hook]&\\
\{\pm1\}\arrow[ur]\arrow[dr]&&\mathcal{G}^+\\
&L_\sigma\arrow[ur,hook]&
\end{tikzcd}}
\]

--- If $\mathcal{G}$ is ``big enough'', we will moreover have that $\{\pm1\}\hookrightarrow\mathcal{G}^+$ i.e. ``$-1$'' $\ne1$ in $\mathcal{G}^+$, and even
\[
L_\rho\cap L_\sigma=\{\pm1\}\ .
\]

The conditions (6) therefore imply that
\[
\tag{10}
\left\{\begin{aligned}
u(L_\rho)&\ \inserted{\mathcal{G}^+\text{-}}\text{conjugate to }L_\rho\\
u(L_\sigma)&\ \mathcal{G}^+\text{-conjugate to }L_\sigma
\end{aligned}\right.
\]

If $\mathcal{G}^+$ is ``big enough'' so that $\rho$ and $\rho^{-1}$, $\sigma$ and $\sigma^{-1}$ are not conjugate in it, then an element $(u,\mu,\varepsilon,\varepsilon')$ of $Z_{\rho,\sigma}$ is known from the sole knowledge of $(u,\mu)$. If on the other hand the homomorphism
\[
\tag{11}
\diagram{\begin{tikzcd}[column sep=large,row sep=small]
\widehat{\mathbb{Z}}\arrow[r]&\mathcal{G}^+\\
n\arrow[r,mapsto]&\varepsilon_0^n
\end{tikzcd}}
\]
is injective, then $\mu$ is known when we know $u$. Thus if $\mathcal{G}^+$ is big enough to satisfy the three conditions considered (for example so that we have a ``modular homomorphism''
\[
\tag{12}
\diagram{\begin{tikzcd}[column sep=large]
\mathcal{G}\arrow[r]&\mathrm{Gl}(2,\widehat{\mathbb{Z}})
\end{tikzcd}}
\]
(sending $\rho,\sigma,\tau'$ to the elements with the same names), then an element of $Z_{\rho,\sigma}$ is known from the sole
\mspage{549}

{\small
\begin{center}
549\\
\texttt{https://grothendieck.umontpellier.fr}\\
Université de Montpellier, UMR CNRS 5149, Case courrier 051, Place Eugène Bataillon\\
34095 Montpellier cedex 5 - France
\end{center}
}

\scanpage{551}

knowledge of $u\in\operatorname{Aut}(\mathcal{G}^+)$.

We have a homomorphism
\[
\tag{13}
\diagram{\begin{tikzcd}[column sep=large,row sep=small]
L_0\arrow[r,"i_{L_0,Z}"]&\unsure{Z_{\rho,\sigma}}\\
\ell\arrow[r,mapsto]&(\operatorname{int}(\ell),1,1,1)
\end{tikzcd}}
\]
whose image is \unsure{normal}, the kernel is none other than $L_0\cap\operatorname{Cent}(\mathcal{G}^+)$. Let us note that we have
\[
\tag{14}
L_0\cap\operatorname{Cent}(\mathcal{G}^+)\subset L_0\cap\operatorname{Cent}_{\mathcal{G}^+}(\rho)\subset L_0\cap L_1
\]
for if $\ell\in L_0$ and $\rho$ commute, \inserted{i.e. \uline{if $\ell$}$=\rho(\ell)$,} since \deleted{$\ell=\rho(\ell)$} $\rho(L_0)=L_1$, we will have $\ell\in L_1$. Thus if $L_0\cap L_1=\{1\}$ (which is the case if (12) exists), \deleted{then} and even if there exists only
\[
\tag{12'}
\diagram{\begin{tikzcd}[column sep=large]
\mathcal{G}\arrow[r]&\mathrm{Gl}(2,\widehat{\mathbb{Z}})'
\end{tikzcd}}\qquad\text{modular},
\]
then we will have
\[
\tag{15}
L_0\cap\operatorname{Cent}(\mathcal{G}^+)=\{1\},
\]
and (13) will be injective. We are going to assume this, to fix ideas.

We put
\[
\tag{16}
\gamma_{\rho,\sigma}=Z_{\rho,\sigma}/\operatorname{Im}L_0,
\]
so that we will have a canonical exact sequence
\[
\tag{17}
\diagram{\begin{tikzcd}[column sep=large]
1\arrow[r]&L_0\arrow[r,"i_{L_0,Z}"]&Z_{\rho,\sigma}\arrow[r,"\delta_Z"]&\gamma_{\rho,\sigma}\arrow[r]&1
\end{tikzcd}}\ .
\]

The canonical projection homomorphisms $(u,\mu,\varepsilon,\varepsilon')\mapsto\mu$, $\varepsilon,\varepsilon'$ factor through $\gamma_{\rho,\sigma}$ and give canonical homomorphisms
\mspage{550}

{\small
\begin{center}
550\\
\texttt{https://grothendieck.umontpellier.fr}\\
Université de Montpellier, UMR CNRS 5149, Case courrier 051, Place Eugène Bataillon\\
34095 Montpellier cedex 5 - France
\end{center}
}

\scanpage{552}
\mspage{579}

\[
\tag{18}
\left\{\diagram{\begin{tikzcd}[column sep=large,row sep=small]
\gamma_{\rho,\sigma}\arrow[r,"\chi_\gamma"]&\widehat{\mathbb{Z}}^*\\
\gamma_{\rho,\sigma}\arrow[r,"\varepsilon_\gamma"]&\{\pm1\}\\
\gamma_{\rho,\sigma}\arrow[r,"\varepsilon'_\gamma"]&\{\pm1\}
\end{tikzcd}}\right.
\quad\text{i.e.}\quad
\diagram{\begin{tikzcd}[column sep=huge]
\gamma_{\rho,\sigma}\arrow[r,"(\chi_\gamma,\varepsilon_\gamma,\varepsilon'_\gamma)"]&\widehat{\mathbb{Z}}^*\times\{\pm1\}\times\{\pm1\}
\end{tikzcd}}
\]
the composites of these with $Z_{\rho,\sigma}\to\gamma_{\rho,\sigma}$ are denoted
\[
\tag{19}
\chi_Z=\chi_\gamma\circ\delta_Z,\quad
\varepsilon_Z=\varepsilon_\gamma\circ\delta_Z,\quad
\varepsilon'_Z=\varepsilon'_\gamma\circ\delta_Z,
\]
so that
\[
\tag{20}
\chi_Z(u,\mu,\varepsilon,\varepsilon')=\mu,\quad
\varepsilon_Z(u,\mu,\varepsilon,\varepsilon')=\varepsilon,\quad
\varepsilon'_Z(u,\mu,\varepsilon,\varepsilon')=\varepsilon'
\]

We would feel like denoting by $\gamma_{\rho,\sigma}^{\unsure{0|\circ}}$ the kernel of $(\varepsilon_\gamma,\varepsilon'_\gamma)$, by $\gamma_{\rho,\sigma}^{+}$ that of $\chi_\gamma$, by $\gamma_{\rho,\sigma}^{\unsure{0|\circ}+}$ their intersection, and introducing under the notations $Z_{\rho,\sigma}^{\unsure{0|\circ}}$, $Z_{\rho,\sigma}^{+}$, $Z_{\rho,\sigma}^{\unsure{0|\circ}+}$ their inverse images in $Z_{\rho,\sigma}$, i.e. the kernels \unsure{resp.$^{\mathrm t}$|resp.} of $(\varepsilon_Z,\varepsilon'_Z)$, of $\chi_Z$ and of $(\chi_Z,\varepsilon_Z,\varepsilon'_Z)$.

\margin{The groups $S_0G_{\rho,\sigma}$, $G_{\rho,\sigma}$}

$3^\circ)$ \deleted{The group $Z_{\rho,\sigma}$ \unsure{operates},} \inserted{We have a tautological homomorphism}
\[
\tag{21}
\diagram{\begin{tikzcd}[column sep=large,row sep=small]
Z_{\rho,\sigma}\arrow[r]&\operatorname{Aut}(\mathcal{G}^+)\\
(u,\mu,\varepsilon,\varepsilon')\arrow[r,mapsto]&u
\end{tikzcd}}\ ,
\]
thus $Z_{\rho,\sigma}$ \unsure{acts|operates} on $\mathcal{G}^+$, and we can therefore consider the $\tfrac12$ direct product
\[
\tag{22}
\begin{gathered}
S_0G_{\rho,\sigma}=Z_{\rho,\sigma}\mathbin{.}\mathcal{G}^+\ .\\[-2pt]
(\tfrac12\text{-direct})
\end{gathered}
\]
\mspage{551}

{\small
\begin{center}
551\\
\texttt{https://grothendieck.umontpellier.fr}\\
Université de Montpellier, UMR CNRS 5149, Case courrier 051, Place Eugène Bataillon\\
34095 Montpellier cedex 5 - France
\end{center}
}

\scanpage{553}

We have a homomorphism
\[
\tag{23}
\diagram{\begin{tikzcd}[column sep=large,row sep=small]
L_0\arrow[r,"i_{L_0,SG}"]&S_0G_{\rho,\sigma}=Z_{\rho,\sigma}\mathbin{.}\mathcal{G}^+\\
\ell\arrow[r,mapsto]&i_{L_0,Z}(\ell)\mathbin{.}\ell^{-1}
\end{tikzcd}}
\]
whose image is again normal, \deleted{\illegible{1}} \inserted{We} \deleted{\unsure{is stable}} \inserted{\unsure{put}}
\[
\tag{24}
G_{\rho,\sigma}=\underbrace{S_0G_{\rho,\sigma}}_{Z_{\rho,\sigma}\mathbin{.}\mathcal{G}^+}/\operatorname{Im}L_0\ \deleted{\unsure{Z_\rho|\mathcal{G}}},
\]
so that we will have an exact sequence
\[
\tag{25}
\diagram{\begin{tikzcd}[column sep=large]
1\arrow[r]&L_0\arrow[r,no head,"i_{L_0,SG}"]&S_0G_{\rho,\sigma}\arrow[r]&G_{\rho,\sigma}\arrow[r,no head]&1
\end{tikzcd}}\ ,
\]
and \unsure{a} diagram \inserted{\unsure{commutative}} of homomorphisms
\[
\tag{26}
\diagram{\begin{tikzcd}[column sep=large,row sep=small]
&Z_{\rho,\sigma}\arrow[dr,hook,"i_Z"]&\\
L_0\arrow[ur,hook,"i_{L_0,Z}"]\arrow[dr,hook,"i_{L_0,\mathcal{G}}"']&&G_{\rho,\sigma}\\
&\mathcal{G}^+\arrow[ur,hook,"i_{\mathcal{G}^+}"']&
\end{tikzcd}}
\]
(a priori $\operatorname{Ker}(\mathcal{G}^+\to G_{\rho,\sigma})\simeq\operatorname{Ker}(L_0\to Z_{\rho,\sigma})\simeq L_0\cap\operatorname{Cent}(\mathcal{G}^+)$, but by hypothesis (15) it is reduced to $\{1\}$), fitting into a commutative diagram of exact sequences
\[
\tag{27}
\diagram{\begin{tikzcd}[column sep=large,row sep=large]
1\arrow[r]&L_0\arrow[r,"i_{L_0,Z}"]\arrow[d,hook,"i_{L_0,\mathcal{G}}"']&Z_{\rho,\sigma}\arrow[r,"\delta_Z"]\arrow[d,hook]&\gamma_{\rho,\sigma}\arrow[r]\arrow[d,equal]&1\\
1\arrow[r]&\mathcal{G}^+\arrow[r,"i_{\mathcal{G}^+}"]&G_{\rho,\sigma}\arrow[r,"\delta_G"]&\gamma_{\rho,\sigma}\arrow[r,no head]&1
\end{tikzcd}}\ .
\]
We will put
\[
\tag{28}
\chi_G=\chi_\gamma\circ\delta_G,\quad
\varepsilon_G=\varepsilon_\gamma\circ\delta_G,\quad
\varepsilon'_\gamma\circ\delta_G
\]
\mspage{552}

{\small
\begin{center}
552\\
\texttt{https://grothendieck.umontpellier.fr}\\
Université de Montpellier, UMR CNRS 5149, Case courrier 051, Place Eugène Bataillon\\
34095 Montpellier cedex 5 - France
\end{center}
}

\scanpage{554}
\mspage{580}

hence a homomorphism
\[
\tag{29}
\diagram{\begin{tikzcd}[column sep=large]
G_{\rho,\sigma} \arrow[r,"{(\chi_G,\varepsilon_G,\varepsilon'_G)}"] &
\widehat{\mathbb Z}^{*}\times\{\pm1\}\times\{\pm1\}
\end{tikzcd}},
\]
leading us to introduce $G_{\rho,\sigma}^{+}$, $G_{\rho,\sigma}^{\unsure{0}}$,
$G_{\rho,\sigma}^{+0}$ as usual.

The injection $i_{\mathcal G^+}:\mathcal G^+\hookrightarrow G_{\rho,\sigma}$
can be extended to an embedding \deleted{\illegible{1}}
(denoted $i_{\mathcal G}$)
\[
\tag{30}
\diagram{\begin{tikzcd}
i_{\mathcal G}:\mathcal G \arrow[r,hook] & G_{\rho,\sigma}
\end{tikzcd}},
\]
which I define by its value on \deleted{$\unsure{\tau'}$}
\inserted{$\tau=\sigma\tau'$}, by
\[
\tag{31}
i_{\mathcal G}(\tau\deleted{'})=
\deleted{\illegible{1}}\,i_Z(\tau_Z)
\]
\margin{If $\tau_\gamma$ denotes the image of $\tau=\tau_G\in G_{\rho,\sigma}$
in $\gamma_{\rho,\sigma}=G_{\rho,\sigma}/\mathcal G^+$,
$\mathcal G$ is identified with the inverse image in $G_{\rho,\sigma}$ of the subgroup
$\{1,\tau_\gamma\}$ of $\gamma_{\rho,\sigma}$.}
\[
\tag{32}
\left\{\begin{aligned}
&\tau_Z\in Z_{\rho,\sigma}\quad\text{defined as}\\
&\tau_Z=\left(
\underbrace{\operatorname{int}(\tau)|\mathcal G^+}_{\in\operatorname{Aut}(\mathcal G^+)},
\underset{\substack{\rotatebox{90}{$\in$}\\\widehat{\mathbb Z}^{*}}}{-1},
\underset{\substack{\rotatebox{90}{$\in$}\\\{\pm1\}}}{-1},
\underset{\substack{\rotatebox{90}{$\in$}\\\{\pm1\}}}{-1}\right).
\end{aligned}\right.
\]

From now on we identify $\mathcal G$ with a subgroup of $G_{\rho,\sigma}$,
and we denote by the same letter elements of $\mathcal G$
(such as $\tau,\tau',\rho,\sigma,\varepsilon_0,\varepsilon_1$, etc)
and their images in $G_{\rho,\sigma}$.
\margin{Note that $\tau'_G\notin Z_{\rho,\sigma}$ on the other hand we have
\deleted{$\unsure{\tau'_G}$} $\tau_G\in Z_{\rho,\sigma}$
\deleted{\illegible{1}}.}
We will moreover have
\[
\tag{32bis}
\chi_G|\mathcal G=\varepsilon_G|\mathcal G=\varepsilon'_G|\mathcal G
=\varepsilon_{\mathcal G}:\left\{\begin{array}{l}
\text{trivial on }\mathcal G^+\\
-1\text{ on }\mathcal G^-=\mathcal G\setminus\mathcal G^+.
\end{array}\right.
\]

\margin{The groups $Z_\rho,Z_\sigma,N_\rho,N_\sigma,SN_\rho,SN_\sigma$\unsure{\dots}}
$4^\circ)$ We \deleted{define} note that the element $\tau'$ of
\deleted{$\unsure{G_{\rho,\sigma}}$} $\mathcal G\subset G_{\rho,\sigma}$
normalizes $L_\rho$ and $L_\sigma$, therefore also normalizes
\[
\tag{33}
\left\{\begin{aligned}
Z_\rho&\overset{\mathrm{dfn}}=\operatorname{Cent}_{G_{\rho,\sigma}}(\rho)
=\operatorname{Cent}_{G_{\rho,\sigma}}(L_\rho),\\
Z_\sigma&\overset{\mathrm{dfn}}=\operatorname{Cent}_{G_{\rho,\sigma}}(\sigma)
=\operatorname{Cent}_{G_{\rho,\sigma}}(L_\sigma)
\end{aligned}\right.
\]
\mspage{553}

{\small
\begin{center}
553\\
\texttt{https://grothendieck.umontpellier.fr}\\
Université de Montpellier, UMR CNRS 5149, Case courrier 051, Place Eugène Bataillon\\
34095 Montpellier cedex 5 - France
\end{center}
}

\scanpage{555}

thus $\{1,\tau'\}$ acts on these groups, and we can consider
\[
\tag{34}
\left\{\begin{aligned}
N_\rho&=\{1,\tau'\}\mathbin{.}Z_\rho,\qquad\deleted{\unsure{N}}\\
N_\sigma&=\{1,\tau'\}\mathbin{.}Z_\sigma
\end{aligned}\right.
\qquad(\text{products }\tfrac12\text{ direct}).
\]

We have canonical exact sequences
\[
\tag{35}
\diagram{\begin{tikzcd}
1\arrow[r]&\unsure{Z_\rho|L_\rho}\arrow[r]&N_\rho\arrow[r,"\varepsilon_\rho"]&\{\pm1\}\arrow[r]&1\\
1\arrow[r]&Z_\sigma\arrow[r]&N_\sigma\arrow[r,"\varepsilon_\sigma"]&\{\pm1\}\arrow[r]&1
\end{tikzcd}},
\]
and a commutative diagram of canonical homomorphisms
\[
\tag{36}
\diagram{\begin{tikzcd}[row sep=small,column sep=large]
&\mathbb D_\rho\arrow[r]&N_\rho\arrow[dr,"i_\rho"]&\\
\{1,\tau'\}\times\{\pm1\}\arrow[ur,hook]\arrow[dr,hook]&&&G_{\rho,\sigma}\\
&\mathbb D_\sigma\arrow[r]&N_\sigma\arrow[ur,"\unsure{i_\sigma}"']&
\end{tikzcd}}
\]
where
\[
\tag{37}
\left\{\begin{aligned}
\mathbb D_\rho&\simeq\mathbb D_6\subset\operatorname{Gl}(2,\mathbb Z)
\quad\text{generated by }\rho,\tau'\\
\mathbb D_\sigma&\simeq\mathbb D_4\subset\operatorname{Gl}(2,\mathbb Z)
\quad\text{\rule{2.7cm}{0.35pt}}\ \sigma,\tau'\\
\{1,\tau'\}\times\{\pm1\}&=\mathbb D_\rho\cap\mathbb D_\sigma
\quad\text{in }\operatorname{Gl}(2,\mathbb Z),
\end{aligned}\right.
\]
\deleted{\unsure{hence} \unsure{the homomorphisms} $\mathbb D_\rho\longrightarrow N_\rho$}
\[
\tag{38}
\left\{\begin{aligned}
\mathbb D_\rho&=\{1,\tau'\}\mathbin{.}L_\rho
\quad(\text{at least if }\rho\text{ of order }6\text{ in }\mathcal G)\\
\mathbb D_\sigma&=\{1,\tau'\}\mathbin{.}L_\sigma
\quad(\text{\rule{1.5cm}{0.35pt}}\ \sigma\text{ of order }4\ \text{\rule{1.4cm}{0.35pt}}),
\end{aligned}\right.
\]
and the homomorphisms
\[
\diagram{\begin{tikzcd}
\mathbb D_\rho\arrow[r]&N_\rho,
\end{tikzcd}}
\qquad
\diagram{\begin{tikzcd}
\mathbb D_\sigma\arrow[r]&N_\sigma
\end{tikzcd}}
\]
are obvious. The homomorphisms $i_\rho,i_\sigma$ too, of course.

Let us recall that
\[
\tag{39}
\left\{\begin{aligned}
\deleted{\unsure{i'_\rho}}i_\rho\text{ injective}
&\Longleftrightarrow\tau'\notin Z_\rho
\Longleftrightarrow\rho^2\ne1
\Longleftrightarrow\rho\ne-1
\Longleftrightarrow\mathcal G^+\text{ is }\unsure{\deleted{\text{isomorphic}}|\text{isomorphic}}
\text{ to a quotient of }\mathbb Z/4\mathbb Z.\\
i_\sigma\text{ injective}
&\Longleftrightarrow\tau'\notin Z_\sigma
\Longleftrightarrow\sigma^2\ne1\quad
\deleted{\unsure{\Longleftrightarrow\sigma\ne\pm1}}
\end{aligned}\right.
\]
\mspage{554}

{\small
\begin{center}
554\\
\texttt{https://grothendieck.umontpellier.fr}\\
Université de Montpellier, UMR CNRS 5149, Case courrier 051, Place Eugène Bataillon\\
34095 Montpellier cedex 5 - France
\end{center}
}

\scanpage{556}
\mspage{581}

Letting $G_{\rho,\sigma}$ act on $\mathcal G^+$, a normal subgroup,
via inner automorphisms, we see that every element of $G_{\rho,\sigma}$ transforms
$L_\rho\subset\mathcal G^+$ into a \inserted{$\mathcal G^+$-}conjugate subgroup,
thus is \unsure{congruent} modulo \deleted{$\unsure{\operatorname{Norm}_{\mathcal G^+}(L_\rho)}$}
$\mathcal G^+$ to an element of
$\operatorname{Norm}_{\unsure{G|\mathcal G}}(L_\rho)=\unsure{\operatorname{Im}}i_\rho(N_\rho)$,
and likewise for $L_\sigma$. In other words, \inserted{setting}
\inserted{\[
\tag{40}
\begin{aligned}
N_\rho^*&\overset{\mathrm{dfn}}=\{x\in N_\rho\mid\varepsilon_\rho(x)=\varepsilon'_G(i_\rho(x))\},\qquad
N_\sigma^*\overset{\mathrm{dfn}}=\{x\in N_\sigma\mid\varepsilon_\sigma(x)=\varepsilon_G(i_\sigma(x))\}
\end{aligned}
\]}
the composite homomorphisms
\[
\tag{\unsure{40}}
\left\{\begin{array}{ll}
\diagram{\begin{tikzcd}[column sep=large]
N_\rho^*\arrow[r,"i_\rho"]&G_{\rho,\sigma}\arrow[r,"\delta_G"]&\gamma_{\rho,\sigma}
\end{tikzcd}},&\delta_\rho\overset{\mathrm{dfn}}=\delta_G\circ i_\rho\\[1.5em]
\diagram{\begin{tikzcd}[column sep=large]
N_\sigma^*\arrow[r,"i_\sigma"]&G_{\rho,\sigma}\arrow[r,"\delta_G"]&\gamma_{\rho,\sigma}
\end{tikzcd}},&\delta_\sigma\overset{\mathrm{dfn}}=\delta_G\circ i_\sigma
\end{array}\right.
\]
epimorphisms.
\deleted{We \unsure{see} \illegible{3}}
\inserted{We thus have \deleted{\illegible{1}}}

\margin{If $g\in G_{\rho,\sigma}$, and
$\varepsilon=\varepsilon_G(g)$, $\varepsilon'=\varepsilon'_G(g)$, then
\[
\begin{aligned}
\operatorname{int}(g)(\rho)&\underset{\mathcal G^+}{\sim}\rho^{\varepsilon'}
=\unsure{\tau^{\nu'}(\rho)},\\
\operatorname{int}(g)(\sigma)&\underset{\mathcal G^+}{\sim}\sigma^\varepsilon
=\unsure{\tau^\nu(\rho)|\tau^\nu(\sigma)}.
\end{aligned}
\]
\deleted{\unsure{thus}}
(where $(-1)^\nu=\varepsilon$, $(-1)^{\nu'}=\varepsilon'$), \inserted{thus}
\deleted{$\exists g\in\unsure{\mathcal G^+}$ such that}
$\exists a=\tau^{\nu'}\mathbin{.}\unsure{a_0|ag}\in N_\rho$
(where $a_0=\unsure{\tau^\nu g}$) such that

a) $g\equiv\unsure{a|a_0}\ (\mathcal G^+)$ and

b) $\varepsilon_\rho(a)=\unsure{\varepsilon'|\varepsilon'_G}
=\unsure{\varepsilon'_G(g)|\varepsilon_G(g)}$ ---

(and idem for $N_\sigma$\unsure{\dots})

thus
\[
\begin{aligned}
\operatorname{Im}\bigl(N_\rho^{\deleted{\unsure{*}}(+1)}
\xrightarrow{\deleted{\unsure{\delta_\rho}}}\gamma_{\rho,\sigma}\bigr)
&\supset\gamma'_{\rho,\sigma}\deleted{\unsure{\varepsilon'}},\\
\operatorname{Im}\bigl(N_\rho^{(-1)}\longrightarrow\gamma_{\rho,\sigma}\bigr)
&\supset\gamma_{\rho,\sigma}\setminus\gamma'_{\rho,\sigma}.
\end{aligned}
\]
But this does \uline{not} prove that we have \unsure{commutativity}!
\unsure{hence probably}
}

\deleted{\unsure{They are not} \illegible{2} (cf. (39))}
\deleted{(\unsure{41}) a) If $\rho^2\ne1$, then we have commutativity in}
\[
\diagram{\deleted{\begin{tikzcd}[column sep=large,row sep=small]
N_\rho\arrow[rr,"\delta_\rho"]
\arrow[dr,"\varepsilon_\rho"']&&
\gamma_{\rho,\sigma}\arrow[dl,"\varepsilon'_\gamma"]\\
\unsure{\varepsilon_\rho}&\{\pm1\}&\\
&\unsure{\{\pm1\}}&\unsure{\{\pm1\}}
\end{tikzcd}}}
\]
\deleted{hence}
\inserted{commutative} diagram of \deleted{$\unsure{\{\pm1\}}$} exact sequences
\[
\tag{41}
\diagram{\begin{tikzcd}[column sep=large]
1\arrow[r]&Z_\rho\arrow[r]\arrow[d,"\text{epimorphism}"]&N_\rho^*\arrow[r,"\varepsilon_\rho"]\arrow[d,"\text{epimorphism}"]&\{\pm1\}\arrow[r]\arrow[d,equal]&1\\
1\arrow[r]&\gamma'_{\rho,\sigma}\arrow[r]&\unsure{\gamma_{\rho,\sigma}|Z_{\rho,\sigma}}\arrow[r,"\varepsilon'_\gamma"]&\{\pm1\}\arrow[r]&1
\end{tikzcd}}
\]
where we have set
\[
\tag{42}
\gamma'_{\rho,\sigma}=\operatorname{Ker}\bigl(
\diagram{\begin{tikzcd}
\gamma_{\rho,\sigma}\arrow[r,"\varepsilon'_\gamma"]&\{\pm1\}
\end{tikzcd}}\bigr);
\]

\deleted{b) If $\sigma^2\ne1$ (cf. (39)) then we have commutativity}
\inserted{Likewise}
\[
\diagram{\deleted{\begin{tikzcd}[column sep=large,row sep=small]
N_\sigma^*\arrow[rr,"\delta_\sigma"]\arrow[dr,"\varepsilon_\sigma"']&&
\gamma_{\rho,\sigma}\arrow[dl,"\varepsilon_\gamma"]\\
&\{\pm1\}&
\end{tikzcd}}},
\]
\unsure{thus|hence this} commutative\ diagram\ of exact sequences
\[
\tag{43}
\diagram{\begin{tikzcd}[column sep=large]
1\arrow[r]&Z_\sigma\arrow[r]\arrow[d,"\text{epimorphism}"]&N_\sigma^*\arrow[r,"\varepsilon_\sigma"]\arrow[d,"\text{epimorphism}"']&\{\pm1\}\arrow[r]\arrow[d,equal]&1\\
1\arrow[r]&\gamma''_{\rho,\sigma}\arrow[r]&\gamma_{\rho,\sigma}\arrow[r,"\varepsilon_\gamma"]&\{\pm1\}\arrow[r]&1
\end{tikzcd}},\qquad\text{where}
\]
\[
\tag{44}
\gamma''_{\rho,\sigma}=\operatorname{Ker}\bigl(
\diagram{\begin{tikzcd}
\gamma_{\rho,\sigma}\arrow[r,"\varepsilon_\gamma"]&\{\pm1\}
\end{tikzcd}}\bigr).
\]
\mspage{555}

{\small
\begin{center}
555\\
\texttt{https://grothendieck.umontpellier.fr}\\
Université de Montpellier, UMR CNRS 5149, Case courrier 051, Place Eugène Bataillon\\
34095 Montpellier cedex 5 - France
\end{center}
}

\scanpage{557}

\uline{N.B.} The notations $\varepsilon_\gamma,\varepsilon'_\gamma:
\gamma_{\rho,\sigma}\longrightarrow\{\pm1\}$ can lead to confusion,
the first \unsure{referring|pertaining} to the action of a $u$ on $L_\sigma$
(not on $L_\rho$). \unsure{A} slightly heavier notation,
but not leading to confusion, would be
$\varepsilon_{\gamma,\rho},\varepsilon_{\gamma,\sigma}$ --- or even
just $\varepsilon_\rho,\varepsilon_\sigma$, the confusion with the homomorphisms of this name
$\unsure{N'_\rho|N_\rho^*}\longrightarrow\{\pm1\}$,
$N_\sigma\longrightarrow\{\pm1\}$ being harmless when
$\rho^2\ne1,\sigma^2\ne1$, because of the compatibilities (41) (43).

We have thus constructed, with the aid \unsure{of|of the} (1) i.e. of $\mathcal G$ equipped with
$\rho,\sigma,\tau'$, groups \inserted{$\mathcal G^+$,}
$G_{\rho,\sigma},\gamma_{\rho,\sigma},Z_{\rho,\sigma},N'_\rho,N'_\sigma$,
fitting into a commutative diagram of groups
\[
\tag{45}
\diagram{\begin{tikzcd}[column sep=large,row sep=small]
&L_0\simeq SZ_{\rho,\sigma}\arrow[rr]\arrow[ddd,hook]&&Z_{\rho,\sigma}
\arrow[ddd,hook,dashed]\arrow[dddrr,dashed,"\delta_Z"]&&&\\
&&SN_\rho^*\arrow[rr]\arrow[ddl,hook]& &N_\rho^*
\arrow[ddl,hook,dashed]\arrow[ddr,dashed,"\delta_\rho"]&&
\widehat{\mathbb Z}^{*}\times\{\pm1\}\times\{\pm1\}\\
&&SN_\sigma^*\arrow[rr]\arrow[dl,hook,dotted]\arrow[dl]&&N_\sigma^*
\arrow[dl,hook,dotted]\arrow[dl]\arrow[dr,dashed,"\delta_\sigma"]&&\\
1\arrow[r]&\mathcal G^+\arrow[rr,hook,"i_{\mathcal G^+}"]\arrow[dr,hook]&&
G_{\rho,\sigma}\arrow[rr,"\delta_G"]&&\gamma_{\rho,\sigma}\arrow[r]\arrow[uur,"{(\chi_\gamma,\varepsilon_{\gamma,\rho},\varepsilon_{\gamma,\sigma})}"]&1\\
&&\mathcal G\arrow[ur,hook]&&&&
\end{tikzcd}}
\]
where we have set
\[
\tag{46}
\left\{\begin{aligned}
SZ_{\rho,\sigma}&=i_{L_0,Z}(L_0)
=\operatorname{Ker}(\delta_Z:Z_{\rho,\sigma}\longrightarrow\gamma_{\rho,\sigma})\\
SN'_\rho&=\operatorname{Ker}(\delta_\rho:N_\rho^*\longrightarrow\gamma_{\rho,\sigma})
=\operatorname{Cent}_{\mathcal G^+}(\rho)\\
SN'_\sigma&=\operatorname{Ker}(\delta_\sigma:N_\sigma^*\longrightarrow\gamma_{\rho,\sigma})
=\operatorname{Cent}_{\mathcal G^+}(\sigma)
\end{aligned}\right.
\]
\margin{N.B. We can thus also denote them by $SZ_\rho,SZ_\sigma$ instead of the
$SN_\rho^*,SN_\sigma^*$.}
\mspage{556}

{\small
\begin{center}
556\\
\texttt{https://grothendieck.umontpellier.fr}\\
Université de Montpellier, UMR CNRS 5149, Case courrier 051, Place Eugène Bataillon\\
34095 Montpellier cedex 5 - France
\end{center}
}

\scanpage{558}
\mspage{582}

\unsure{in} such a way that we have exact sequences
\[
\tag{47}
\left\{
\diagram{\begin{tikzcd}[column sep=large]
1\arrow[r]&S\unsure{\deleted{Z}N}_\rho^*\arrow[r]&N_\rho^*\arrow[r,"\delta_\rho"]&
\gamma_{\rho,\sigma}\arrow[r]&1\\
1\arrow[r]&SN_\sigma^*\arrow[r]&N_\sigma^*\arrow[r,"\delta_\sigma"]&
\gamma_{\rho,\sigma}\arrow[r]&1
\end{tikzcd}}
\right.
\]
\margin{We have
\[
\tag{47bis}
\left\{\begin{aligned}
\operatorname{Im}\bigl(\unsure{SN_\rho}\longrightarrow\mathcal G^+\bigr)
&=\operatorname{Norm}_{\mathcal G^+}(L_\rho)\\
\operatorname{Im}\bigl(SN_\sigma\longrightarrow\mathcal G^+\bigr)
&=\operatorname{Norm}_{\mathcal G^+}(L_\sigma)
\end{aligned}\right.
\]
(or again, if $\rho^2\ne1$ resp. $\sigma^2\ne1$,
$N_\rho\cap\mathcal G$ and $N_\sigma\cap\mathcal G$)}
mapping to the exact sequence
\[
\diagram{\begin{tikzcd}[column sep=large]
1\arrow[r]&\mathcal G^+\arrow[r]&G_{\rho,\sigma}\arrow[r]&
\gamma_{\rho,\sigma}\arrow[r]&1
\end{tikzcd}},
\]
in a way analogous to (27). Moreover, the homomorphisms
$i_\rho:N_\rho^*\longrightarrow G_{\rho,\sigma}$ and
$i_\sigma:N_\sigma^*\longrightarrow G_{\rho,\sigma}$ have
a marked tendency to be injective (cf. (39)),
which we have marked on diagram (45) by the sign
$\diagram{\begin{tikzcd}[column sep=1.7em]
\phantom{0}\arrow[r,hook,dotted]\arrow[r]&\phantom{0}
\end{tikzcd}}$
(dotted inclusion sign).

We also introduce
\[
\tag{48}
\left\{\begin{aligned}
\chi_\rho&=\chi_\gamma\circ\delta_\rho:
\diagram{\begin{tikzcd}
\unsure{\deleted{Z}N}_\rho^*\arrow[r]&\widehat{\mathbb Z}^{*}
\end{tikzcd}}\\
\chi_\sigma&=\chi_\gamma\circ\unsure{\delta_{\deleted{\rho}\sigma}}:
\diagram{\begin{tikzcd}
\unsure{\deleted{Z}N}_\sigma^*\arrow[r]&\widehat{\mathbb Z}^{*}
\end{tikzcd}}
\end{aligned}\right.,
\]
whose kernels (which contain $S\unsure{\deleted{Z}N}_\rho\deleted{\illegible{1}}$
resp. $SN_\sigma$, and which we \unsure{shall take care not} to confuse with them)
are denoted as is only right by $N_\rho^{*+},N_\sigma^{*+}$.

\uline{N.B.} We should still examine the composites
\[
\tag{49}
\left\{
\diagram{\begin{tikzcd}[column sep=large]
N_\rho^*\arrow[r]&\gamma_{\rho,\sigma}\arrow[r,"\unsure{\varepsilon_{\gamma,\sigma}}"]&\{\pm1\}\\
N_\sigma^*\arrow[r]&\gamma_{\rho,\sigma}\arrow[r,"\varepsilon_{\gamma,\rho}"]&\{\pm1\}
\end{tikzcd}}
\right.,
\]
we shall see that in the ``good cases'', the restrictions of these homomorphisms to
$Z_\rho,Z_\sigma$ respectively are epimorphisms.
\mspage{557}

{\small
\begin{center}
557\\
\texttt{https://grothendieck.umontpellier.fr}\\
Université de Montpellier, UMR CNRS 5149, Case courrier 051, Place Eugène Bataillon\\
34095 Montpellier cedex 5 - France
\end{center}
}

\scanpage{559}
\mspage{\unsure{583}}

\deleted{We set}
\[
\deleted{
\left\{
\begin{aligned}
S\unsure{Z_\rho|N_\rho}
 &=\operatorname{Ker}\bigl(\unsure{Z_\rho|N_\rho}
       \xrightarrow{\delta_\rho}\gamma_{\rho,\sigma}\bigr)
       =\operatorname{Norm}_{\mathcal G}(L_\rho)\\
SN_\sigma
 &=\operatorname{Ker}\bigl(N_\sigma
       \xrightarrow{\delta_\sigma}\gamma_{\rho,\sigma}\bigr)
       =\operatorname{Norm}_{\mathcal G}(L_\sigma)
\end{aligned}
\right.
}
\tag{15}
\]
\[
\deleted{
S\unsure{B}_{\rho,\sigma}
=\operatorname{Ker}\bigl(\unsure{B}_{\rho,\sigma}
   \xrightarrow{\delta_B}\gamma_{\rho,\sigma}\bigr)
\simeq L_0\subset
\left\{
\begin{gathered}
\unsure{\operatorname{Norm}|\operatorname{Cent}}_{\mathcal G}
       (\unsure{L_0|L_\rho})\\[-2pt]
\cap\\[-2pt]
\operatorname{Norm}_{\mathcal G}\,\unsure{L_0|L_\sigma}
\end{gathered}
\right.
}
\]
\margin{Canceled\par
\deleted{\unsure{It seems that}
\unsure{in the case that} \unsure{we want to consider},
\(\operatorname{Norm}_{\mathcal G}(L_0)
\unsure{\ne|=}\unsure{\operatorname{Cent}}_{\mathcal G}(L_0)\)
and idem for \(L_\rho,L_\sigma\),}}

\textcircled{\(5^\circ\)}
\margin{The groups
\(S_0\mathcal M_{\rho,\sigma},\mathcal M_{\rho,\sigma},
S_0\Gamma_{\rho,\sigma},\Gamma_{\rho,\sigma}\).}

We set
\[
\left\{
\begin{aligned}
S_0\mathcal M_{\rho,\sigma}
 &=\deleted{\unsure{B}}Z_{\rho,\sigma}
   \mathop{\times}_{\gamma_{\rho,\sigma}}N_\rho^*
   \mathop{\times}_{\gamma_{\rho,\sigma}}N_\sigma^*
   \mathop{\times}_{\gamma_{\rho,\sigma}}G_{\rho,\sigma}.\\
\mathcal M_{\rho,\sigma}
 &=N_\rho^*
   \mathop{\times}_{\gamma_{\rho,\sigma}}N_\sigma^*
   \mathop{\times}_{\gamma_{\rho,\sigma}}G_{\rho,\sigma}
\end{aligned}
\right.
\tag{50}
\]
\[
\deleted{\unsure{42}}
\qquad
\left\{
\begin{aligned}
S_0\Gamma_{\rho,\sigma}
 &=Z_{\rho,\sigma}
   \mathop{\times}_{\gamma_{\rho,\sigma}}N_\rho^*
   \mathop{\times}_{\gamma_{\rho,\sigma}}N_\sigma^*\\
\Gamma_{\rho,\sigma}
 &=\unsure{N_\rho^*}
   \mathop{\times}_{\gamma_{\rho,\sigma}}N_\sigma^*,
\end{aligned}
\right.
\tag{51}
\]

whence \deleted{\unsure{homomorphism}}\inserted{\uline{diagram}} of exact sequences
\diagram{
\((52)\qquad\)
\begin{tikzcd}[column sep=large,row sep=large]
 & & 1 \arrow[d] & 1 \arrow[d] & \\
 & & L_0 \arrow[d] \arrow[r,equal] & L_0 \arrow[d] & \\
1 \arrow[r] & \mathcal G^+ \arrow[r] \arrow[d,equal]
 & S_0\mathcal M_{\rho,\sigma} \arrow[r] \arrow[d]
 & S_0\Gamma_{\rho,\sigma} \arrow[r] \arrow[d] & 1\\
1 \arrow[r] & \mathcal G^+ \arrow[r]
 & \mathcal M_{\rho,\sigma} \arrow[r] \arrow[d]
 & \Gamma_{\rho,\sigma} \arrow[r] \arrow[d] & 1,\\
 & & 1 & 1 &
\end{tikzcd}
}

and exact sequences
\diagram{
\((53)\qquad\)
\begin{tikzcd}[column sep=large]
1 \arrow[r]
 & \underset{\simeq L_0}{\uline{SZ_{\rho,\sigma}}}
   \times SN_\rho^*\times SN_\sigma^*\times\mathcal G^+
       \arrow[r]
 & S_0\mathcal M_{\rho,\sigma}\arrow[r]
 & \gamma_{\rho,\sigma}\arrow[r] & 1
\end{tikzcd}
}
\diagram{
\((54)\qquad\)
\begin{tikzcd}[column sep=large]
1 \arrow[r]
 & SN_\rho^*\times SN_\sigma^*\times\mathcal G^+\arrow[r]
 & \mathcal M_{\rho,\sigma}\arrow[r]
 & \gamma_{\rho,\sigma}\arrow[r] & 1
\end{tikzcd}
}
\diagram{
\((55)\qquad\)
\begin{tikzcd}[column sep=large]
1 \arrow[r]
 & \underset{\simeq L_0}{\uline{SZ_{\rho,\sigma}}}
   \times SN_\rho^*\times SN_\sigma^*\arrow[r]
 & S_0\Gamma_{\rho,\sigma}\arrow[r]
 & \gamma_{\rho,\sigma}\arrow[r] & 1
\end{tikzcd}
}
\diagram{
\((56)\qquad\)
\begin{tikzcd}[column sep=large]
1 \arrow[r]
 & SN_\rho^*\times SN_\sigma^*\arrow[r]
 & \Gamma_{\rho,\sigma}\arrow[r]
 & \gamma_{\rho,\sigma}\arrow[r] & 1.
\end{tikzcd}
}

\mspage{558}

{\small
\begin{center}
558\\
\texttt{https://grothendieck.umontpellier.fr}\\
Université de Montpellier, UMR CNRS 5149, Case courrier 051, Place Eugène Bataillon\\
34095 Montpellier cedex 5 - France
\end{center}
}

\scanpage{560}

We also have
\[
S_0\mathcal M_{\rho,\sigma}
\simeq\mathcal M_{\rho,\sigma}
       \mathop{\times}_{\gamma_{\rho,\sigma}}Z_{\rho,\sigma}
\tag{57}
\]
\[
S_0\Gamma_{\rho,\sigma}
\simeq\Gamma_{\rho,\sigma}
       \mathop{\times}_{\gamma_{\rho,\sigma}}Z_{\rho,\sigma}
\tag{58}
\]

i.e. the extensions \(S_0\mathcal M_{\rho,\sigma},S_0\Gamma_{\rho,\sigma}\)
of \(\mathcal M_{\rho,\sigma},\Gamma_{\rho,\sigma}\) by \(L_0\)
are deduced from the extension \(Z_{\rho,\sigma}\) of
\(\gamma_{\rho,\sigma}\) by \(L_0\) by inverse image under
\(\mathcal M_{\rho,\sigma}\to\gamma_{\rho,\sigma},
\Gamma_{\rho,\sigma}\to\gamma_{\rho,\sigma}\) --- and we also have
\deleted{\unsure{$\Gamma$}}
\[
S_0\mathcal M_{\rho,\sigma}
\simeq\mathcal M_{\rho,\sigma}
       \mathop{\times}_{\Gamma_{\rho,\sigma}}S_0\Gamma_{\rho,\sigma}
\tag{59}
\]

i.e. the extension \(S_0\mathcal M_{\rho,\sigma}\) of
\(\mathcal M_{\rho,\sigma}\) by \(L_0\) is deduced from the extension
\(S_0\Gamma_{\rho,\sigma}\) of \(\Gamma_{\rho,\sigma}\) by \(L_0\),
by inverse image \unsure{under}
\(\mathcal M_{\rho,\sigma}\to\Gamma_{\rho,\sigma}\).

\inserted{We introduce an element}
\[
\deleted{\unsure{\tau'_{\mathcal M}}},
\qquad
\tau'_{\deleted{\unsure{S_0}}\mathcal M}
  \in\deleted{S_0}\mathcal M_{\rho,\sigma}
\quad
\deleted{\text{\unsure{the}}}
\quad
\tau'_{\deleted{\unsure{S_0}}\mathcal M}
 =\deleted{\unsure{\tau'_Z}}
 \bigl(
  \overset{\inserted{\substack{a\\\Vert}}}{\tau'_\rho},
  \overset{\inserted{\substack{b\\\Vert}}}{\tau'_\sigma},
  \overset{\inserted{\substack{u\\\Vert}}}{\tau'_G}
 \bigr)
\tag{60}
\]

\deleted{\unsure{of homomorphisms} between the groups
introduced, I am going to link them in
\deleted{in}\inserted{a} commutative diagram}
\inserted{where \(\deleted{\unsure{\tau'_Z}},\tau'_\rho,\tau'_\sigma,
\tau'_G\) are the}
``avatars'' of \(\tau'\), taken resp. in
\(\deleted{\unsure{Z_{\rho,\sigma}}}
N_\rho^*,N_\sigma^*,G_{\rho,\sigma}\).
These do indeed all have the same image in
\(G_{\rho,\sigma}\) and a fortiori in \(\gamma_{\rho,\sigma}\),
hence define an \unsure{element} of
\(\deleted{\unsure{G}}\inserted{\mathcal M}_{\rho,\sigma}\),
whose operation on \(\mathcal G^+\) is that of \(\tau\).
Whence an embedding
\diagram{
\((61)\qquad\)
\begin{tikzcd}
\mathcal G \arrow[r,hook,"i_{\mathcal G,\mathcal M}"]
 & \mathcal M_{\rho,\sigma}
\end{tikzcd}
\quad : identity on \(\mathcal G^+\)
\qquad
\begin{tikzcd}
\tau'_{\mathcal G}\arrow[r,mapsto] & \tau'_{\mathcal M}.
\end{tikzcd}
}

\mspage{559}

{\small
\begin{center}
559\\
\texttt{https://grothendieck.umontpellier.fr}\\
Université de Montpellier, UMR CNRS 5149, Case courrier 051, Place Eugène Bataillon\\
34095 Montpellier cedex 5 - France
\end{center}
}

\scanpage{561}
\mspage{584}

\diagram{
\((62)\qquad\)
\begin{tikzcd}[column sep=small,row sep=large]
 & & L_0\times SN_\rho^*\times SN_\sigma^*
 & & & & L_0 & & & &\\
 & & & & S_0\Gamma_{\rho,\sigma}
 & & & & Z_{\rho,\sigma} & &\\
L_0\times SN_\rho^*\times SN_\sigma^*\times\mathcal G^+
 & & SN_\rho^*\times SN_\sigma^*
 & & & & \mathrm{III}^{\unsure{+|'}}
 & & & &\\
 & & & S_0\mathcal M_{\rho,\sigma}
 & \Gamma_{\rho,\sigma} & & N_\rho^*
 & & & & \{\pm1\}\\
 & SN_\rho^*\times SN_\sigma^*\times\mathcal G^+
 & & & & & N_\sigma^*
 & & \gamma_{\rho,\sigma} & & \widehat{\mathbb Z}^{*}\\
 & & & & & \mathrm{II}
 & \deleted{\unsure{\mathcal G^+|N_\sigma^*}}
 & & & & \{\pm1\}\\
 & & & \mathcal M_{\rho,\sigma}
 & & & & & G_{\rho,\sigma} & &\\
\mathcal G^+ & & & & & & & & & &\\
 & & \mathcal G & & & & & & & &\\
 & & \Vert & & & & & & & &\\
 & & \{1,\tau'\}\mathbin{.}\mathcal G^+
 & & & & & & & &
\arrow[from=1-3,to=1-7,"\mathrm{pr}"]
\arrow[from=1-3,to=2-5,hook]
\arrow[from=1-3,to=3-3,"\mathrm{pr}"]
\arrow[from=1-7,to=2-9,hook]
\arrow[from=2-5,to=2-9,"\theta_Z"]
\arrow[from=2-5,to=4-5]
\arrow[from=2-9,to=5-9,"\delta_Z"]
\arrow[from=3-1,to=1-3,"\mathrm{pr}"]
\arrow[from=3-1,to=5-2,"\mathrm{pr}"']
\arrow[from=3-1,to=4-4,hook]
\arrow[from=3-3,to=4-5,hook]
\arrow[from=4-4,to=2-5]
\arrow[from=4-4,to=7-4]
\arrow[from=4-5,to=4-7,"\theta_\rho"]
\arrow[from=4-5,to=5-7,"\theta_\sigma"']
\arrow[from=4-7,to=5-9,"\delta_\rho"]
\arrow[from=5-7,to=5-9,"\delta_\sigma"]
\arrow[from=5-2,to=3-3,"\mathrm{pr}"']
\arrow[from=5-2,to=7-4,hook]
\arrow[from=5-2,to=6-7,dashed,"\unsure{\mathrm{pr}}"]
\arrow[from=6-7,to=7-9,dashed,hook]
\arrow[from=7-4,to=4-5]
\arrow[from=7-4,to=6-7,dotted]
\arrow[from=7-4,to=7-9,"\theta_G"]
\arrow[from=7-9,to=5-9,"\delta_G"']
\arrow[from=5-9,to=4-11,
 "{\varepsilon_{\gamma,\rho}=\varepsilon'_\gamma}"]
\arrow[from=5-9,to=5-11,"\chi_\gamma"]
\arrow[from=5-9,to=6-11,
 "{\varepsilon_{\gamma,\sigma}=\varepsilon_\gamma}"']
\arrow[from=8-1,to=5-2,hook,"\uline{\mathrm{inj}}"]
\arrow[from=8-1,to=7-4,dashed]
\arrow[from=8-1,to=9-3,hook]
\arrow[from=9-3,to=7-4,hook,"i_{\mathcal G,\mathcal M}"]
\arrow[from=9-3,to=7-9,dashed,hook,"i_{\mathcal G,G}"']
\arrow[from=1-3,to=2-9,phantom,
 "{\deleted{\unsure{\mathrm{IV}}}\ \mathrm{V}}" description,pos=.52]
\arrow[from=4-7,to=5-7,phantom,"\mathrm{I}" description]
\arrow[from=8-1,to=7-4,phantom,"\mathrm{IV}" description,pos=.5]
\end{tikzcd}
}

We have summarized \deleted{\unsure{the main}}\inserted{some of the}
relations between the groups introduced so far in the commutative
diagram (62), which I comment on briefly.

a) All the arrows are either injections \((\hookrightarrow)\)
or surjections \((\longrightarrow)\), with the exception
\inserted{perhaps} of the three arrows issuing from \(\gamma_{\rho,\sigma}\).

b) The ten squares (six faces of a cube, and
\deleted{four}\inserted{\unsure{three|the}} squares marked
\(\mathrm I,\mathrm{II},\mathrm{III}^{\unsure{+|'}},
\mathrm{IV},\mathrm V\)) are cartesian.

c) The \inserted{\unsure{five}} arrows marked \(\mathrm{pr}\)
are projections, canonical from product groups onto partial
products. The arrow marked \(\mathrm{inj}\)
\diagram{
\begin{tikzcd}
\mathcal G^+\arrow[r] &
SN_\rho^*\times SN_\sigma^*\times\mathcal G^+
\end{tikzcd}
}
is the canonical injection.

d) In order not to destroy the commutativity of the diagram,
we have omitted to show in it the arrows
\inserted{(of inclusion, in particular)}
\diagram{
\begin{tikzcd}
Z_{\rho,\sigma}\arrow[r] & G_{\rho,\sigma},
\end{tikzcd}
\qquad
\begin{tikzcd}
N_\rho^*\arrow[r] & G_{\rho,\sigma},
\end{tikzcd}
\qquad
\begin{tikzcd}
N_\sigma^*\arrow[r] & G_{\rho,\sigma}
\end{tikzcd}
}
which appeared in (45).

\mspage{560}

{\small
\begin{center}
560\\
\texttt{https://grothendieck.umontpellier.fr}\\
Université de Montpellier, UMR CNRS 5149, Case courrier 051, Place Eugène Bataillon\\
34095 Montpellier cedex 5 - France
\end{center}
}

\scanpage{562}

e) In order not to overload the diagram, I have not introduced in
the diagram notations for certain useful composite arrows,
such as
\[
(63)\quad
\left\{
\begin{array}{ll}
\delta_\Gamma:\Gamma_{\rho,\sigma}\longrightarrow\gamma_{\rho,\sigma},
&
\deleted{\unsure{\delta_{S_0\Gamma}}:
\unsure{S_0\Gamma_{\rho,\sigma}}\longrightarrow
\unsure{\gamma_{\rho,\sigma}}.}\\
\delta_{S_0\Gamma}:S_0\Gamma_{\rho,\sigma}
\longrightarrow\gamma_{\rho,\sigma}&\\
\delta_{\mathcal M}:\mathcal M_{\rho,\sigma}
\longrightarrow\gamma_{\rho,\sigma}&\\
\delta_{S_0\mathcal M}:S_0\mathcal M_{\rho,\sigma}
\longrightarrow\gamma_{\rho,\sigma}&
\end{array}
\right.
\]
and
\[
(64)\quad
\left\{
\begin{aligned}
\chi_\Gamma&=\chi_\gamma\circ\delta_\Gamma:
\Gamma_{\rho,\sigma}\longrightarrow\widehat{\mathbb Z}^{*}\\
\chi_{S_0\Gamma}&=\chi_\gamma\circ\delta_{S_0\Gamma}:
S_0\Gamma_{\rho,\sigma}\longrightarrow\widehat{\mathbb Z}^{*}\\
\chi_{\mathcal M}&=\chi_\gamma\circ\delta_{\mathcal M}:
\mathcal M_{\rho,\sigma}\longrightarrow\widehat{\mathbb Z}^{*}\\
\chi_{S_0\mathcal M}&=\chi_\gamma\circ\delta_{S_0\mathcal M}:
S_0\mathcal M_{\rho,\sigma}\longrightarrow\widehat{\mathbb Z}^{*}
\end{aligned}
\right.
\]
and the composites
\(\varepsilon_{?,\rho},\varepsilon_{?,\sigma}:?\longrightarrow\{\pm1\}\),
where \(?\) can take \deleted{\unsure{any}} value
\(Z,G\) (already introduced) \(N_\rho^*,N_\sigma^*\)
(we will have
\(\varepsilon_{N_\rho,\rho}=\varepsilon_\rho
\deleted{\text{\unsure{simply}}},\quad
\varepsilon_{N_\sigma,\sigma}=\varepsilon_\sigma
\deleted{\text{\unsure{simply}}}\)),
\(\Gamma,S_0\Gamma,\mathcal M,S_0\mathcal M\).
The kernel of \(\chi_?\) is denoted \(?^+\), that of
\(\varepsilon_{?,\rho}\) is denoted \(?'\), that of
\(\varepsilon_{?,\sigma}\) is denoted \(?''\),
we write
\[
?^\circ\doteq?'\cap?'',\qquad
?^{+'}=?^+\cap?',\quad
?^{+''}=?^+\cap?'',
\]
\[
?^{+,\circ}=?^+\cap?^\circ,\quad
S?=\operatorname{Ker}\delta_?,
\]
\deleted{and} hence \(\deleted{\unsure{S?\subset}}\)

\margin{\deleted{there is a small incompatibility with the notation
\(N'_\rho,N'_\sigma\) already used in another \unsure{sense}.}}

\[
(65)
\]
\diagram{
\begin{tikzcd}[row sep=1.15em,column sep=1.6em]
 & & ?^{'+} & &\\
\deleted{\unsure{S?\subset}}\quad ?^{\circ+}
 & & & & ?^+\ \deleted{\text{\unsure{we}}}\\
 & & ?^{''+} & &\\
 & & \unsure{?'} & &\\
?^\circ & & & & ?\\
 & & ?'' & &
\arrow[from=2-1,to=1-3,hook]
\arrow[from=1-3,to=2-5,hook]
\arrow[from=2-1,to=3-3,hook]
\arrow[from=3-3,to=2-5,hook]
\arrow[from=2-1,to=5-1,hook]
\arrow[from=1-3,to=4-3,hook,shift left=.7ex]
\arrow[from=3-3,to=6-3,hook,shift right=.7ex]
\arrow[from=2-5,to=5-5,hook]
\arrow[from=5-1,to=4-3,hook]
\arrow[from=5-1,to=6-3,hook]
\arrow[from=4-3,to=5-5,hook]
\arrow[from=6-3,to=5-5,hook]
\end{tikzcd}
}
``cartesian'' cube.

f) Let us also recall that we have commutativity in
\[
(66)
\]
\diagram{
\begin{tikzcd}[row sep=.8em,column sep=2em]
 & \{\pm1\}\arrow[dr] &\\
\mathcal G
\arrow[ur,"{\varepsilon_{\mathcal G,\rho}}"]
\arrow[rr,"\chi_{\mathcal G}"]
\arrow[dr,"{\varepsilon_{\mathcal G,\sigma}}"']
 & & \widehat{\mathbb Z}^{*}\\
 & \{\pm1\}\arrow[ur] &
\end{tikzcd}
}
, i.e. \(\varepsilon_{\mathcal G}:\mathcal G\longrightarrow\{\pm1\}\),

\mspage{561}

{\small
\begin{center}
561\\
\texttt{https://grothendieck.umontpellier.fr}\\
Université de Montpellier, UMR CNRS 5149, Case courrier 051, Place Eugène Bataillon\\
34095 Montpellier cedex 5 - France
\end{center}
}

\scanpage{563}
\mspage{585}

The kernel of $\varepsilon_{\mathcal{G}}$ is none other than $\mathcal{G}^{+}$. In \unsure{its capacity} as a subgroup of $\mathcal M_{\rho,\sigma}$ (resp. of $G_{\rho,\sigma}$) $\mathcal{G}$ is identified with the inverse image in this group of the subgroup $\{1,\tau'_{\Gamma}\}$ of $\Gamma_{\rho,\sigma}$ (resp. \inserted{of the subgroup} $\{1,\tau'_{\gamma}\}$ of $\gamma_{\rho,\sigma}$), where $\tau'_{\Gamma},\tau'_{\gamma}$ denote the images of $\tau'\in\mathcal{G}$ in $\Gamma_{\rho,\sigma}$ and $\gamma_{\rho,\sigma}$ respectively.

g) The diagram (62) has as its main part
\[
\text{(67)}\qquad
\diagram{\begin{tikzcd}[row sep=0.9em,column sep=1.25em]
{} & {}\arrow[rrrr]\arrow[ddr] & {} & {} & {} & {}\arrow[dd] & {} & {} & {} & {}\\
{} & {} & {} & {}\arrow[drr] & {} & {} & {} & {} & {} & {}\\
{}\arrow[uur]\arrow[ddr] & {} & {}\arrow[ur]\arrow[dr] & {} & {} & {}\arrow[urrr,dashed,shorten <=3em]\arrow[rrr,shorten <=3em]\arrow[drrr,dashed,shorten <=3em] & \Biggl[ & {} & {} & \Biggr]\\
{} & {} & {} & {}\arrow[urr] & {} & {} & {} & {} & {} & {}\\
\deleted{\unsure{[\mathcal{G}^{+}]}} & {}\arrow[uur]\arrow[rrrr] & {} & {} & {} & {}\arrow[uu] & {} & {} & {} & {}\\
\unsure{\Bigl[}{}\arrow[ur,dashed,hook]\arrow[urrrrr,dashed,hook] & {} & {} & {} & {} & \Bigr] & {} & {} & {} & {}
\end{tikzcd}},
\]
all the groups marked on this \inserted{last} diagram, except for $\widehat{\mathbb Z}^{*},\{\pm1\},\{\pm1\}$ (put in brackets $[\;]$) map to $\gamma_{\rho,\sigma}$, which is the terminal term of this diagram, linking \unsure{9} groups (not \unsure{counting} $\mathcal{G}\ \deleted{\unsure{\mathcal{G}^{+}}}$ inserted \deleted{\unsure{by}} as \deleted{\unsure{terms}}\inserted{\unsure{outside}} \deleted{\unsure{initial} of}\inserted{of two} dotted arrows) including $\gamma_{\rho,\sigma}$, obtained from four of them $Z_{\rho,\sigma},N_{\rho}^{*},N_{\sigma}^{*},G_{\rho,\sigma}$ (which map to $\gamma_{\rho,\sigma}$ by epimorphisms) by taking fiber products over $\gamma_{\rho,\sigma}$ of certain of them --- \deleted{\unsure{the}} the initial term being $S_{0}\mathcal M_{\rho,\sigma}$, which is the fiber product of the four preceding factors. Taking the kernels of the homomorphisms from these \deleted{four} nine groups to $\gamma_{\rho,\sigma}$, we find a diagram of the same combinatorial type, linking

\mspage{562}

{\small
\begin{center}
562\\
\texttt{https://grothendieck.umontpellier.fr}\\
Université de Montpellier, UMR CNRS 5149, Case courrier 051, Place Eugène Bataillon\\
34095 Montpellier cedex 5 - France
\end{center}
}

\scanpage{564}

nine groups, \deleted{including $\{1\}=\operatorname{Ker}(\gamma_{\rho,\sigma}\longrightarrow\gamma_{\rho,\sigma})$,} formed from $L_{0}=SZ_{\rho,\sigma}$, $SN_{\rho}^{*},SN_{\sigma}^{*}$ and $\mathcal{G}^{+}=SG_{\rho,\sigma}$ and from certain products of these groups (including the product over the empty index set, namely $\{1\}=\operatorname{Ker}(\gamma_{\rho,\sigma}\xrightarrow{\mathrm{id}}\gamma_{\rho,\sigma})$), with as \deleted{\unsure{arrowisms|arrows}} canonical projections. We have inserted in (62) only \deleted{five} of these product groups, with the canonical injections into the fiber products from which they arise, \deleted{in} dispensing with including in it \unsure{also|also} $SN_{\rho}^{*},SN_{\sigma}^{*},SG_{\rho,\sigma}=\mathcal{G}^{+}$, and the canonical inclusions $SN_{\rho}^{*}\hookrightarrow N_{\rho}^{*}$, $SN_{\sigma}^{*}\hookrightarrow N_{\sigma}^{*}$, $\mathcal{G}^{+}=SG_{\rho,\sigma}\hookrightarrow G_{\rho,\sigma}$, and the \deleted{\unsure{p}} \inserted{complete} diagram of the projections of the partial products of these groups and of $L_{0}$, \deleted{including} \inserted{which therefore \unsure{lacks}} the projections onto these \inserted{three} groups, and the homomorphisms from these groups and from $L_{0}$ to the unit group.

\mspage{563}

{\small
\begin{center}
563\\
\texttt{https://grothendieck.umontpellier.fr}\\
Université de Montpellier, UMR CNRS 5149, Case courrier 051, Place Eugène Bataillon\\
34095 Montpellier cedex 5 - France
\end{center}
}

\scanpage{565}
\mspage{586}

6°) \uline{The set} $\Sigma_{\rho,\sigma},S\unsure{\mathcal{G}|g}_{\rho,\sigma}$.

We put
\[
\text{(68)}\qquad
\begin{gathered}
\Sigma_{\rho,\sigma}\subset\mathcal{G}^{+}\times\mathcal{G}^{+}\times\widehat{\mathbb Z}^{*}\times\{\pm1\}\times\{\pm1\}\\
\begin{array}{c}
\Vert\\[-0.4em]
\displaystyle
\left\{(\xi,\eta,\mu,\unsure{\varepsilon',\varepsilon})\;\middle|\;
\begin{array}{l}
\xi\rho\ \inserted{\mathcal{G}^{+}\text{-}}\text{conjugate to }\rho^{\varepsilon'}\\
\eta\sigma\ \mathcal{G}^{+}\text{-conjugate to }\sigma^{\varepsilon}\\
\xi=\eta\underbrace{\varepsilon_{1}^{\mu-1}}_{\rho_{1}^{\nu}\ \text{where }\nu=\frac{\mu-1}{2}}\quad\deleted{\unsure{\text{(with}}}
\end{array}\right\}
\end{array}
\end{gathered}
\]

\uline{N.B.} If $\rho$ not conjugate to $\rho^{-1}$ in $\mathcal{G}^{+}$, $\varepsilon'$ known when we know $\xi,\eta$ --- likewise if $\sigma$ not conjugate to $\sigma^{-1}$ (hence $\sigma^{2}\ne1$), $\varepsilon'$ is determined by $\xi,\eta$. \unsure{If this homomorphism} $\widehat{\mathbb Z}\longrightarrow L_{0}$ $(n\longmapsto\varepsilon_{0}^{n})$ is injective, then $\mu$ is determined by $\xi,\eta$, and the relation $\xi=\eta\varepsilon_{1}^{\mu-1}$ can also be written (without bringing in $\mu$)
\[
\text{(69)}\qquad \eta^{-1}\xi\in L_{1}.
\]

Let on the other hand
\[
\text{(70)}\qquad
\begin{gathered}
S\mathcal{G}_{\rho,\sigma}\subset\mathcal{G}\times\mathcal{G}\times\widehat{\mathbb Z}^{*}\\
\begin{array}{c}
\Vert\\[-0.4em]
\displaystyle\left\{(\alpha,\beta,\mu)\;\middle|\;
\underbrace{[\alpha,\rho]}_{\xi}=[\beta,\sigma]\varepsilon_{1}^{\mu-1}\right\}
\end{array}
\end{gathered}
\]

We have a map
\[
\text{(71)}\qquad
\left\{\diagram{\begin{tikzcd}[row sep=1.5em,column sep=3em]
S\mathcal{G}_{\rho,\sigma}\arrow[r] & \Sigma_{\rho,\sigma}\\
(\alpha,\beta,\mu)\arrow[r,mapsto] & (\xi,\eta,\mu,\varepsilon',\varepsilon)
\end{tikzcd}}\right.
\quad\text{with}\quad
\left\{\begin{array}{l}
\xi=[\alpha,\rho]\\
\eta=[\beta,\sigma]\\
\mu\ \text{the same}\\
\varepsilon'=\unsure{\varepsilon}_{\mathcal{G}}(\alpha)\\
\varepsilon=\varepsilon_{\mathcal{G}}(\beta)
\end{array}\right.
\]

\margin{\uline{N.B.} $\xi=\alpha\rho(\alpha)^{-1}\in\mathcal{G}^{+}$, since we have $\varepsilon_{\mathcal{G}}(\xi)=\unsure{\dot{\alpha}|\bar{\alpha}}\rho(\unsure{\dot{\alpha}|\bar{\alpha}})=1$ since $\rho$ operates trivially on $\mathcal{G}/\mathcal{G}^{+}\simeq\{\pm1\}$, for the same reason $\eta=\beta\sigma(\beta)^{\unsure{-1}}\in\mathcal{G}^{+}$.}

which is \inserted{obviously} \uline{surjective}.

We define on the other hand

\mspage{564}

{\small
\begin{center}
564\\
\texttt{https://grothendieck.umontpellier.fr}\\
Université de Montpellier, UMR CNRS 5149, Case courrier 051, Place Eugène Bataillon\\
34095 Montpellier cedex 5 - France
\end{center}
}

\scanpage{566}

\[
\text{(72)}\qquad
\begin{gathered}
\Sigma_{\rho,\sigma}^{\mathrm{eff}}\subset\Sigma_{\rho,\sigma}\\
\begin{array}{c}
\Vert\\[-0.4em]
\displaystyle\left\{(\xi,\unsure{\eta},\mu,\varepsilon',\varepsilon)\in\Sigma_{\rho,\sigma}\;\middle|\;
\begin{array}{l}
\exists u\in\operatorname{Aut}(\mathcal{G}^{+})\ \text{with}\\
u(\rho)=\xi\rho,\quad u(\unsure{\sigma})=\eta\sigma
\end{array}\right\}
\end{array}
\end{gathered}
\]
\margin{i.e. $[u,\rho]=\xi$, $[u,\sigma]=\eta$}
(this $u$ will then be unique, \inserted{and \unsure{conversely}, \unsure{knowledge} of $u$ determines $\xi,\eta$.}) There is a $1$-$1$ correspondence between the elements of $\Sigma_{\rho,\sigma}^{\mathrm{eff}}$, and the systems $(u,\mu,\varepsilon',\varepsilon)\in\operatorname{Aut}(\mathcal{G}^{+})\times\widehat{\mathbb Z}^{*}\times\{\pm1\}\times\{\pm1\}$, such that we have \deleted{$u(\rho)\ \mathcal{G}^{+}\text{-conjugate}\unsure{\text{.}}$}
\[
\left\{\begin{array}{lll}
u(\rho)\text{ is}&\mathcal{G}^{+}\text{-conjugate to}&\rho^{\varepsilon'}\\
u(\sigma)&\rule{9em}{0.35pt}&\sigma^{\varepsilon}\\
u(\varepsilon_{0})=\varepsilon_{0}^{\mu}&\multicolumn{2}{l}{(\text{which expresses }\xi=\eta\rho_{1}^{\nu})}
\end{array}\right.
\]
--- now this set is none other than the subgroup $\unsure{Z_{\rho,\sigma}|\mathbb Z_{\rho,\sigma}}$ of $\operatorname{Aut}(\mathcal{G}^{+})\times\widehat{\mathbb Z}^{*}\times\{\pm1\}\times\{\pm1\}$. \inserted{(cf. (6) page\,577')}

We thus find a bijection
\[
\text{(73)}\qquad
\diagram{\begin{tikzcd}[row sep=2em,column sep=3em]
Z_{\rho,\sigma}\arrow[r,"\sim"] & \Sigma_{\rho,\sigma}^{\mathrm{eff}}\\
\begin{array}{c}\deleted{\unsure{(u,\mu,\varepsilon,\varepsilon')}}\\u\end{array}\arrow[r,mapsto] & (\xi,\eta,\mu,\varepsilon',\varepsilon)
\end{tikzcd}}
\quad\left\{\begin{array}{l}
\xi=[u,\rho]\\
\eta=[u,\sigma]\\
\mu=\chi(u)\\
\varepsilon'=\varepsilon'_{Z}(u)\\
\varepsilon=\varepsilon_{Z}(u)
\end{array}\right.\ .
\]

We denote by $S\mathcal{G}_{\rho,\sigma}^{\mathrm{eff}}$ the inverse image of $\Sigma_{\rho,\sigma}^{\mathrm{eff}}$ in $S\mathcal{G}_{\rho,\sigma}$ under (71), so we will have
\[
\text{(74)}\qquad
\begin{gathered}
S\mathcal{G}_{\rho,\sigma}^{\mathrm{eff}}\subset S\mathcal{G}_{\rho,\sigma}\\
\begin{array}{c}
\Vert\\[-0.4em]
\displaystyle\left\{\unsure{\{\alpha,\beta,\mu\}}\in\mathcal{G}\times\mathcal{G}\times\widehat{\mathbb Z}^{*}\;\middle|\;
\begin{array}{l}
[\alpha,\rho]=[\beta,\sigma]\varepsilon_{1}^{\mu-1}\\
\exists u\in\operatorname{Aut}(\mathcal{G}^{+})\ \text{such that we have}\\
\left\{\begin{array}{ll}
u(\rho)=\alpha(\rho)&(=\xi\rho,\text{ where }\xi=[\alpha,\rho])\\
u(\sigma)=\beta(\sigma)&(=\eta\sigma,\text{ where }\eta=[\beta,\sigma])
\end{array}\right.
\end{array}\right\}
\end{array}
\end{gathered}
\]
so that we have a \uline{surjective} \unsure{homomorphism} induced by (71)
\[
\text{(75)}\qquad
\left|\diagram{\begin{tikzcd}[column sep=3em]
S\mathcal{G}_{\rho,\sigma}^{\mathrm{eff}}\arrow[r]&\Sigma_{\rho,\sigma}^{\mathrm{eff}}
\end{tikzcd}}\right.\ .
\]

\mspage{565}

{\small
\begin{center}
565\\
\texttt{https://grothendieck.umontpellier.fr}\\
Université de Montpellier, UMR CNRS 5149, Case courrier 051, Place Eugène Bataillon\\
34095 Montpellier cedex 5 - France
\end{center}
}

\scanpage{567}
\mspage{587}

We have a canonical map
\[
\text{(76)}\qquad
\left\{\diagram{\begin{tikzcd}[row sep=2em,column sep=2em]
S_{0}\Gamma_{\rho,\sigma}=\unsure{Z|\mathbb Z}_{\rho,\sigma}\times_{\gamma_{\rho,\sigma}}N_{\rho}^{*}\times_{\gamma_{\rho,\sigma}}N_{\sigma}^{*}\arrow[r]&S\mathcal{G}_{\rho,\sigma}^{\mathrm{eff}}\\
x=\unsure{\{u,a,b\}|(u,a,b)}\arrow[r,mapsto] & (\alpha,\beta,\mu)
\end{tikzcd}}\right.
\quad\text{with}\quad
\left\{\begin{array}{l}
\alpha=u\overbrace{a^{-1}\tau'^{\unsure{\nu'(x)}}}^{\in Z_{\rho}}\\
\beta=u\underbrace{b^{-1}\tau'^{\nu(x)}}_{\in Z_{\sigma}}\\
\mu=\chi_{S\Gamma}(x)
\end{array}\right.
\]

\margin{to be more correct, one should write
\[\left\{\begin{array}{l}
\alpha=u\,i_{\rho}(a)^{-1}\tau'^{\nu'(x)}\\
\beta=u\,i_{\sigma}(b)^{-1}\tau'^{\nu(x)}
\end{array}\right.\]}

where we have defined $\nu'(x),\unsure{\nu(x)|\nu'(x)}\in\mathbb Z/2\mathbb Z$ by
\[
\deleted{\unsure{(}}\qquad
(-1)^{\nu'(x)}=\varepsilon'(x),\qquad(-1)^{\nu(x)}=\varepsilon(x).
\]

By definition of $S_{0}\Gamma_{\rho,\sigma}$ as a fiber product, we will indeed have $ua^{-1},ub^{-1}\in\mathcal{G}^{+}$, hence $\alpha,\beta\in\mathcal{G}$, we will moreover have
\[
\text{(77)}\qquad
\left\{\begin{array}{l}
\displaystyle\alpha(\rho)=u\left(\underbrace{a^{-1}\left(\underbrace{\tau'^{\nu'(x)}(\rho)}_{\rho^{\varepsilon'(x)}}\right)}_{(\rho^{\varepsilon'(x)})^{\varepsilon'(x)}=\rho}\right)=u(\rho),\quad\text{and likewise}\\[2em]
\beta(\sigma)=u(\sigma),
\end{array}\right.
\qquad\text{i.e.}\quad
\begin{array}{l}
[\alpha,\rho]=[u,\rho]\\
[\beta,\sigma]=[u,\sigma]
\end{array}
\]

\margin{\uline{N.B.} We have, for \unsure{every} $a\in N_{\rho}^{*}$, if $\varepsilon_{\rho}(a)=(-1)^{\nu_{\rho}(a)}$, $a\tau'^{\nu_{\rho}(a)}\in Z_{\rho}$, hence $[ua\tau'^{\nu_{\rho}(a)},\rho]=[u,\rho]$. Likewise for $\sigma\ldots$}

\deleted{so that} moreover the condition
$\underbrace{[\alpha,\rho]}_{[u,\rho]}=\underbrace{[\beta,\sigma]}_{[u,\sigma]}\varepsilon_{1}^{\mu-1}$ \deleted{\unsure{here}}
merely expresses the condition that $u(\varepsilon_{0})=\varepsilon_{0}^{\mu}$ which enters into the very definition of $Z_{\rho,\sigma}$.

\deleted{Moreover} \inserted{I say that} the map in question \inserted{(76)} is injective, \deleted{for} It suffices to see that when $(\alpha,\beta,\mu)$ are known, we know $u$ --- \inserted{indeed} we \deleted{have \unsure{then}} \inserted{in terms of $(\alpha,\beta,\mu)$}
\[
\text{(78)}\qquad
\left\{\begin{array}{l}
\nu'(x)=\unsure{\nu|\varepsilon}_{\mathcal{G}}(\alpha),\quad
\nu(x)=\unsure{\nu|\varepsilon}_{\mathcal{G}}(\beta),\quad
\deleted{\unsure{\mu=}}\ \chi_{S\Gamma}(x)=\unsure{\chi}_{Z}(u)=\mu\\[0.6em]
\deleted{\text{whence we obtain}}\quad\text{whence}\quad
a=\tau'^{\unsure{\nu|\varepsilon}_{\mathcal{G}}(\alpha)}\alpha^{-1}u,\quad
b=\tau'^{\unsure{\nu|\varepsilon}_{\mathcal{G}}(\beta)}\beta^{-1}u
\end{array}\right.
\]
\deleted{\unsure{we}} so if $u$ is known, we also know $a$ and $b$ hence $x$. On the other hand to determine \inserted{in terms of $\alpha,\beta,\mu$,} $u\in Z_{\rho,\sigma}\subset\operatorname{Aut}(\mathcal{G}^{+})\times\widehat{\mathbb Z}^{*}\times\{\pm1\}\times\{\pm1\}$ it suffices, in addition to (78), to determine \deleted{\unsure{it}} \deleted{\unsure{as}} the element of $\operatorname{Aut}(\mathcal{G}^{+})$ which it defines, which is

\mspage{566}

{\small
\begin{center}
566\\
\texttt{https://grothendieck.umontpellier.fr}\\
Université de Montpellier, UMR CNRS 5149, Case courrier 051, Place Eugène Bataillon\\
34095 Montpellier cedex 5 - France
\end{center}
}
